\originalchapter{规范对称性、Higgs场与宇宙磁单极问题}
本章对应8.286第21--22讲，特别采用第22讲Scott A. Hughes代课逐字稿，并补齐其示意讨论中的数学条件。原课程教师为Alan H. Guth，但不能把本讲口述全部署为Guth。QFT卷给出规范理论的一般结构，本章只解释热宇宙中对称性变化为何关系到物质质量与遗迹；电弱实验机制和假设的大统一相变分别说明。

\originalsection{从电磁势的自由度到局部相位}
取$\hbar=c=1$，电荷为$q$的场满足$\psi'=\ee^{\ii q\Lambda(x)}\psi$。若$D_\mu=\partial_\mu-\ii q A_\mu$且$A'_\mu=A_\mu+\partial_\mu\Lambda$，直接展开得
\[
 D'_\mu\psi'=\ee^{\ii q\Lambda}D_\mu\psi.
\]
所以由$|D\psi|^2$建立的动力学不因任意局部相位改变。$F_{\mu\nu}=\partial_\mu A_\nu-\partial_\nu A_\mu$也保持不变。这把经典势的规范自由度与量子场局部相位联系起来；改变描述不会改变电场、磁场或能量密度。

$U(1)$由单位模复数$\ee^{\ii\theta}$组成，它有一个连续参数。$SU(n)$由$n\times n$酉矩阵且行列式为1组成，连续参数数目$n^2-1$。\emph{$SU(1)$只有单位元素，不能代替$U(1)$}。非Abel群的生成元不对易，导致规范场强含$[A_\mu,A_\nu]$并允许规范玻色子自相互作用；这本身不让玻色子获得质量。没有破缺的QCD仍有无质量的拉格朗日胶子，而Abel规范场也可通过Higgs机制获得质量。原代课口述将质量与非Abel性联系过紧，本卷按动力学区别重写。

标准模型局部规范结构是$SU(3)_c\times SU(2)_L\times U(1)_Y$，局部规范场共$8+3+1$种。弱$SU(2)$在内部空间作用，不等同于真实三维空间旋转；$SU(2)$是$SO(3)$的双覆盖而不是同一个群。Higgs的“双重态”指弱同位旋的两个复分量，两分量都是Lorentz标量，不能称为自旋$1/2$的时空旋量。

\originalsection{电弱Higgs机制中的势与质量}
设Higgs弱双重态$H=(H_1,H_2)^T$，超荷$Y=1/2$。最简单势为
\[
 V(H)=-\mu_H^2 H^\dagger H+\lambda(H^\dagger H)^2+V_0,
 \quad \mu_H^2,\lambda>0.
\]
令$x=H^\dagger H$，驻点$-\mu_H^2+2\lambda x=0$给$x=v^2/2$，其中$v^2=\mu_H^2/\lambda$。可选规范使最低能背景$H=(0,v/\sqrt2)^T$，并将径向涨落写为$(0,(v+h)/\sqrt2)^T$。二阶势给$m_h^2=2\lambda v^2$。

协变导数$D_\mu=\partial_\mu-\ii g\tau^aW^a_\mu/2-\ii g'YB_\mu$，将背景代入$|DH|^2$，得到
\[
 -\mathcal L_{\rm mass}=
 \frac{g^2v^2}{4}W^+_\mu W^{-\mu}
 +\frac{v^2}{8}(gW^3_\mu-g'B_\mu)^2,
\]
这里使用负时间号度规，质量拉格朗日项为所列二次型的负值；质量平方由系数确定。故$m_W=gv/2$、$m_Z=v\sqrt{g^2+g'^2}/2$；正交组合$A_\mu=(g'W^3_\mu+gB_\mu)/\sqrt{g^2+g'^2}$无质量。残余电磁生成元是$Q=T_3+Y$，对所选背景作用为零。准确结构是$SU(2)_L\times U(1)_Y$到$U(1)_{\rm em}$，不能把超荷$U(1)$直接说成未混合的光子。

带Yukawa耦合的费米子有$m_f=y_fv/\sqrt2$，其相对论能谱的静止能不再为零。电荷量子数由群表示规定，不是在Higgs场冷却时从无到有随意创造；左手电子与中微子同处一个弱双重态，也不表示所有场和相互作用完全相同。最小标准模型没有右手中微子场，通常的可重整化Higgs耦合因而不给活跃中微子质量；振荡说明需要额外结构。本卷不把特定质量机制当成已经实验选定。

\originalsection{热相变与大统一假设的区别}
高温有效势会含热修正，例如$c_TT^2H^\dagger H$。当正热项大于零温负二次项时，最低点可接近零；降温后非零径向背景成为有利状态。这是一个解释结构的简化势，电弱转变的实际阶数与模型参数、涨落和规范处理有关，不能由图中一条曲线认定真实宇宙必然先形成尖锐泡壁。

第21--23讲讨论$SU(5)$大统一。其生成元数目24，标准模型可嵌入其中；超荷方向可写成与
\[
 \operatorname{diag}(-1/3,-1/3,-1/3,1/2,1/2)
\]
成比例的无迹矩阵。这展示颜色三维块、弱二维块和一个Abel方向怎样共存于更大的群，不是大统一已存在的证明。最小$SU(5)$的高能破缺通常涉及与电弱双重态不同的Higgs表示。规范耦合随能标的演化可能提示统一；超对称粒子会改变演化斜率，但是否存在相应粒子和高能理论须独立检验。课程$10^{16}$ GeV指假设的大统一尺度，不能误写成已确立电弱Higgs转变的温度。

\originalsection{什么样的真空空间允许拓扑缺陷}
若一套场的势最低点组成空间$\mathcal M$，不同区域可选择不同最低能配置。场往往能连续调整以消除差异，只有不能连续收缩的边界配置才强迫留下稳定缺陷。离散不连通最低点可给畴壁；绕圆一周不能缩掉的环路可给弦；包围点缺陷的二球映射若不能缩为一点，可给单极，分别由$\pi_0(\mathcal M),\pi_1(\mathcal M),\pi_2(\mathcal M)$描述。

这些拓扑条件决定是否\emph{容许}缺陷，生成丰度还需相变动力学。两个区域随意有不同Higgs相位，并不自动证明其中夹着单极。电弱Higgs双重态最低点的情形也不能无条件套用大统一单极结论。课程示意中的畴壁、宇宙弦与磁单极是不同破缺模型的可能结果，不是所有转变的统一三步演化过程。

一个可看清单极的例子是$SU(2)$规范场与实伴随三重态$\Phi^a$，势最低点$\Phi^a\Phi^a=v_G^2$。远离核心的“刺猬”边界$\Phi^a\simeq v_Gx^a/r$把空间二球一对一映到内部二球，绕数为1；若场处处连续且从不离开最低点，就无法将此边界延伸到实心球中心，故核心需偏离最低势。这个绕数条件才是非平凡结的原因。

单独标量的角方向梯度$\partial\Phi\sim v_G/r$，远处能量$\int r^2|\partial\Phi|^2\dd r\sim v_G^2\int\dd r$发散。$1/r$是\emph{角向}梯度，不意味着场的模长必随$\ln r$增长。规范场可使$D_i\Phi$在远处趋零，并留下一条$B_r\propto r^{-2}$的残余磁场；其远区能量$\int r^2B^2\dd r$收敛。完整有限核心轮廓与非Abel单极解需解耦合场方程，这里给出拓扑和有限能量的入口，不声称已证明一般GUT的全部解。

\originalsection{磁荷与质量的量级}
Gaussian单位下，若最小相关电荷为$q$，Dirac量子化给$qg_m=n\hbar c/2$。取$q=e,n=1$作为课程约定，则
\[
 \frac{g_m}{e}=\frac1{2\alpha_{\rm em}}\simeq68.5.
\]
两粒子在同一距离下磁Coulomb力的大小比单位电荷力约大$68.5^2$；异号吸引，同号排斥。含分数电荷、具体群的整体结构和磁荷谱时须重新检查最小允许$n$，不能不加条件地把所有GUT单极都等同于这一最小Dirac荷。

典型平滑GUT单极的质量能量量级为$M_Xc^2/\alpha_G$，其中$M_X$是相关重规范玻色子的质量，$\alpha_G$是高能规范耦合；它不一定等于今日的电磁精细结构常数。课程以$10^{16}$ GeV和约$10^2$倍作粗估，得$10^{18}$ GeV的质量能量，相当于约$1.8$微克。它强调“一个粒子也可宏观重”的可能性，不是已测到这种粒子的质量。

\originalsection{过多遗迹如何成为宇宙学问题}
在非暴胀的热相变里，场的关联长度$\xi$不能远超过当时因果尺度。若按每关联体积约一个单极估计，$n_M\sim\xi^{-3}$；取最大的$\xi\sim c/H$给粗略的低丰度估计。自然单位下，辐射时期$H=1.66\sqrt{g_*}T^2/M_P$、$s=2\pi^2g_{*s}T^3/45$，故
\[
 Y_M=\frac{n_M}{s}\sim10.4\frac{g_*^{3/2}}{g_{*s}}
 \left(\frac T{M_P}\right)^3.
\]
忽略后续湮灭且保持熵史时，$\Omega_Mh^2\simeq2.74\times10^8(M_Mc^2/\mathrm{GeV})Y_M$。假设$T=10^{16}$ GeV、$g_*=g_{*s}=100$和$M_Mc^2=10^{18}$ GeV，量级远大于1，约$10^{19}$。具体数值随模型和关联长度改变；原课$10^{20}$是数量级警示，并非精密丰度预测。

单极形成后若经历$N$个e折，数密度被稀释$\ee^{-3N}$；场的势能却可近恒定，并在结束时转为新的热浴。它解释暴胀为何可移走旧遗迹又保留普通物质。不过再加热若再次超过同一相变的产生阈值，问题可能重现，需规定模型热历史和残余湮灭。下面将这一约束与暴胀的结束机制合起来讨论。
